\documentclass[12pt]{article}
\usepackage[utf8]{inputenc}
\usepackage[brazil]{babel}
\usepackage[T1]{fontenc}

\begin{document}

\section{Metas do Mês de Abril}
As atividades planejadas e executadas durante o mês de abril foram:

\begin{enumerate}
    \item Compreender o funcionamento do framework Flutter;
    \item Configurar o ambiente de desenvolvimento local;
    \item Criar contas no GitHub e capacitar a equipe em comandos básicos de Git.
\end{enumerate}

\subsection{Funcionamento do Flutter}
A equipe realizou pesquisas sobre a arquitetura e a aplicabilidade do Flutter. Os resultados estão sintetizados abaixo:

O Flutter é um framework de código aberto criado pelo Google, amplamente utilizado no mercado por sua versatilidade. Seu principal diferencial é o desenvolvimento **multiplataforma**, permitindo a criação de aplicativos para Android, iOS, Web e Desktop utilizando um único código-base. Isso torna o ciclo de desenvolvimento mais ágil, eficiente e economicamente viável.

No que tange à interface do usuário (UI), o Flutter utiliza uma abordagem declarativa onde a interface é construída inteiramente através de \textbf{Widgets}. Widgets são componentes estruturais e visuais pré-definidos e altamente personalizáveis. O objetivo do projeto é garantir que a interface seja interativa, intuitiva e acessível, proporcionando uma experiência de uso dinâmica.

\subsection{Configuração do Ambiente de Desenvolvimento}
Foi realizada a instalação local do SDK do Flutter e da linguagem Dart. Para garantir a paridade entre os ambientes de desenvolvimento de todos os membros, as versões das dependências foram gerenciadas e travadas no arquivo \texttt{pubspec.yaml} e no \texttt{pubspec.lock}. Uma reunião técnica será agendada para padronizar o ambiente nas máquinas de todos os colaboradores.

\subsection{GitHub e Controle de Versão}
Todos os membros do projeto já possuem contas ativas no GitHub. A próxima etapa consiste no treinamento da equipe quanto ao fluxo de trabalho com Git, utilizando exemplos práticos dentro do repositório oficial do projeto para fixação dos comandos de sincronização.
Consegui dar uma passada geral no básico de GitHub para o pessoal do grupo (\texttt{pull}, \texttt{commit}, \texttt{push}).
Em pouco tempo teremos a capacitação de GitHub do ramo, assim, essa etapa será 100\% completa.

\section{Metas do Mês de Maio}
As atividades planejadas para este período focam na transição do conhecimento teórico para a implementação prática das funcionalidades core do aplicativo:

\begin{enumerate}
    \item Domínio dos fundamentos da linguagem Dart;
    \item Desenvolvimento de interfaces (telas) básicas em Flutter;
    \item Implementação das funcionalidades primárias do ecossistema do aplicativo;
    \item Estruturação e divisão técnica de tarefas entre os membros.
\end{enumerate}

\subsection{Fundamentos de Dart}
No curso disponilibilizado pelo veterano, existe um módulo totalmente focado em Dart, assim, conseguimos assimilar o conteúdo, e concluir essa meta

\subsection{Desenvolvimento de Telas em Flutter}
Realizamos uma reunião entre nós para começar o desenvolvimento, Porém a equipe enfrentou dificuldades, por isso, fizemos uma reunião com um veterano experiente em Flutter. Ele nos deu uma boa direção, nos ajudando a pensar melhor nas funcionalidades, e disponibilizando um curso Produzido por ele mesmo.

\subsection{Desenvolvimento de Funcionalidades e Repositório}
O repositório oficial no GitHub foi configurado e todos os colaboradores foram devidamente integrados como contribuidores. A estrutura base do projeto \textit{skills\_up} foi estabelecida, com o foco de desenvolvimento voltado para a apresentação na \textbf{Feira das Profissões}. As funcionalidades prioritárias definidas são:
\begin{itemize}
    \item \textbf{Lista de Tarefas (To-do list):} Gerenciamento dinâmico de atividades;
    \item \textbf{Temporizador (Pomodoro):} Ferramenta para gestão de tempo e foco;
    \item \textbf{Calendário:} Organização e visualização de prazos e eventos.
\end{itemize}

\subsection{Divisão de Tarefas}
Para otimizar o fluxo de trabalho e garantir a entrega de um Produto Mínimo Viável (MVP) de alta qualidade, a equipe foi segmentada em frentes de trabalho específicas, conforme a tabela abaixo:

\begin{itemize}
    \item \textbf{Design de Interfaces:} Arthur Vilela e Jefferson Noriyuki;
    \item \textbf{Módulo To-do List:} Rafael Teixeira e Lucas Che Kai;
    \item \textbf{Módulo Timer (Pomodoro):} Marcello Buccelli e Miguel de Souza;
    \item \textbf{Módulo Calendário:} Arthur Romanha e Bruno de Lima.
\end{itemize}

\end{document}
